\documentclass[10pt,a4paper]{amsart}
\usepackage{amsmath,amssymb,amsthm,graphicx}
\usepackage{verbatim}
\newenvironment{reference}{\comment}{\endcomment}
%\newenvironment{reference}{}{}
\newenvironment{slogan}{\comment}{\endcomment}
\newenvironment{history}{\comment}{\endcomment}

% For commutative diagrams we use Xy-pic
\usepackage[all]{xy}

% We use 2cell for 2-commutative diagrams.
\xyoption{2cell}
% Theorem environments.
%
\theoremstyle{plain}
\newtheorem{theorem}[subsection]{Theorem}
\newtheorem{proposition}[subsection]{Proposition}
\newtheorem{lemma}[subsection]{Lemma}

\theoremstyle{definition}
\newtheorem{definition}[subsection]{Definition}
\newtheorem{example}[subsection]{Example}
\newtheorem{exercise}[subsection]{Exercise}
\newtheorem{situation}[subsection]{Situation}

\theoremstyle{remark}
\newtheorem{remark}[subsection]{Remark}
\newtheorem{remarks}[subsection]{Remarks}

\numberwithin{equation}{subsection}

% Macros
%
\def\lim{\mathop{\mathrm{lim}}\nolimits}
\def\colim{\mathop{\mathrm{colim}}\nolimits}
\def\Spec{\mathop{\mathrm{Spec}}}
\def\Hom{\mathop{\mathrm{Hom}}\nolimits}
\def\Ext{\mathop{\mathrm{Ext}}\nolimits}
\def\SheafHom{\mathop{\mathcal{H}\!\mathit{om}}\nolimits}
\def\SheafExt{\mathop{\mathcal{E}\!\mathit{xt}}\nolimits}
\def\Sch{\mathit{Sch}}
\def\Mor{\mathop{\mathrm{Mor}}\nolimits}
\def\Ob{\mathop{\mathrm{Ob}}\nolimits}
\def\Sh{\mathop{\mathit{Sh}}\nolimits}
\def\NL{\mathop{N\!L}\nolimits}
\def\CH{\mathop{\mathrm{CH}}\nolimits}
\def\proetale{{pro\text{-}\acute{e}tale}}
\def\etale{{\acute{e}tale}}
\def\QCoh{\mathit{QCoh}}
\def\Ker{\mathop{\mathrm{Ker}}}
\def\Im{\mathop{\mathrm{Im}}}
\def\Coker{\mathop{\mathrm{Coker}}}
\def\Coim{\mathop{\mathrm{Coim}}}

% Boxtimes
%
\DeclareMathSymbol{\boxtimes}{\mathbin}{AMSa}{"02}

%
% Macros for moduli stacks/spaces
%
\def\QCohstack{\mathcal{QC}\!\mathit{oh}}
\def\Cohstack{\mathcal{C}\!\mathit{oh}}
\def\Spacesstack{\mathcal{S}\!\mathit{paces}}
\def\Quotfunctor{\mathrm{Quot}}
\def\Hilbfunctor{\mathrm{Hilb}}
\def\Curvesstack{\mathcal{C}\!\mathit{urves}}
\def\Polarizedstack{\mathcal{P}\!\mathit{olarized}}
\def\Complexesstack{\mathcal{C}\!\mathit{omplexes}}
% \Pic is the operator that assigns to X its picard group, usage \Pic(X)
% \Picardstack_{X/B} denotes the Picard stack of X over B
% \Picardfunctor_{X/B} denotes the Picard functor of X over B
\def\Pic{\mathop{\mathrm{Pic}}\nolimits}
\def\Picardstack{\mathcal{P}\!\mathit{ic}}
\def\Picardfunctor{\mathrm{Pic}}
\def\Deformationcategory{\mathcal{D}\!\mathit{ef}}
\usepackage{cleveref}
\crefname{originalreference}{original reference}{original references}
\Crefname{originalreference}{Original reference}{Original references}
\makeatletter
\newlabel{spaces-morphisms-lemma-locally-quasi-finite-separated-representable}{{0418}{source}{}{sourcepointer-1}{}}
\newlabel{spaces-morphisms-lemma-locally-quasi-finite-separated-representable@cref}{{[originalreference][0][]0418}{[1][0][]source}{}{}{}}
\makeatother
\crefname{originalreference}{original reference}{original references}
\Crefname{originalreference}{Original reference}{Original references}
\makeatletter
\newlabel{spaces-morphisms-lemma-monomorphism-loc-finite-type-loc-quasi-finite}{{0463}{source}{}{sourcepointer-2}{}}
\newlabel{spaces-morphisms-lemma-monomorphism-loc-finite-type-loc-quasi-finite@cref}{{[originalreference][0][]0463}{[1][0][]source}{}{}{}}
\makeatother
\crefname{originallemma}{original lemma}{original lemmas}
\Crefname{originallemma}{Original lemma}{Original lemmas}
\makeatletter
\newlabel{lemma-base-change-invariance-of-relative-dimension}{{0DRN}{source}{}{sourcepointer-3}{}}
\newlabel{lemma-base-change-invariance-of-relative-dimension@cref}{{[originallemma][0][]0DRN}{[1][0][]source}{}{}{}}
\makeatother
\crefname{originallemma}{original lemma}{original lemmas}
\Crefname{originallemma}{Original lemma}{Original lemmas}
\makeatletter
\newlabel{lemma-irreducible-implies-equidimensional}{{0DRX}{source}{}{sourcepointer-4}{}}
\newlabel{lemma-irreducible-implies-equidimensional@cref}{{[originallemma][0][]0DRX}{[1][0][]source}{}{}{}}
\makeatother
\crefname{originalreference}{original reference}{original references}
\Crefname{originalreference}{Original reference}{Original references}
\makeatletter
\newlabel{etale-lemma-finite-unramified-etale-local}{{04HJ}{source}{}{sourcepointer-5}{}}
\newlabel{etale-lemma-finite-unramified-etale-local@cref}{{[originalreference][0][]04HJ}{[1][0][]source}{}{}{}}
\makeatother
\crefname{originalreference}{original reference}{original references}
\Crefname{originalreference}{Original reference}{Original references}
\makeatletter
\newlabel{descent-lemma-dimension-at-point-local}{{04N4}{source}{}{sourcepointer-6}{}}
\newlabel{descent-lemma-dimension-at-point-local@cref}{{[originalreference][0][]04N4}{[1][0][]source}{}{}{}}
\makeatother
\title{A finite type monomorphism from an irreducible stack of dimension d into a reduced stack of dimension <= d is generically an open immersion onto an irreducible component}
\author{The Stacks Project Authors}
\date{stacks-project:stacks-geometry.tex:local-snapshot-20261012; GNU FDL 1.2 or later}
\begin{document}
\maketitle
\noindent\textit{Editorial note.} All delivered mathematical text is the author's own, copied verbatim from the pinned Stacks Project chapter file stacks-geometry.tex; no statement or proof step was rewritten, reordered or removed.\par\medskip
\noindent\textit{Editorial note.} Span recomputation by label pairing (LANE\_BRIEF 10.2): in this chapter the author's proof environment follows the statement environment as a standalone block; the spec covers the complete environment pair (statement span, proof span) as recorded in the pool candidate.\par\medskip
\noindent\textit{Editorial note.} Occupation scan: combined claim\_key + statement\_anchor + generic-line scan over all registration files (active and withdrawn) found 0 hits for stacks-geometry.tex (ge-occupied.json, 2026-10-14); 207 route=geometry / chapter-list field hits in focused-seeds are alias noise (route name and alphabetical chapter enumerations), excluded per the ladder alias-blind-spot caveat.\par\medskip
\ifdefined\hypertarget\hypertarget{sourcepointer-1}{}\fi
\noindent\textit{Source reference.} Original reference 0418: sources/stacks-chow/tags.txt, source lines 5283--5283. Cross-chapter standard prerequisite cited by the author inside the counted span; it is used as a pointer and its own proof is not part of this item.\par\medskip
\ifdefined\hypertarget\hypertarget{sourcepointer-2}{}\fi
\noindent\textit{Source reference.} Original reference 0463: sources/stacks-chow/tags.txt, source lines 5458--5458. Cross-chapter standard prerequisite cited by the author inside the counted span; it is used as a pointer and its own proof is not part of this item.\par\medskip
\ifdefined\hypertarget\hypertarget{sourcepointer-3}{}\fi
\noindent\textit{Source reference.} Original lemma 0DRN: sources/stacks-chow/tags.txt, source lines 17228--17228. Cross-chapter standard prerequisite cited by the author inside the counted span; it is used as a pointer and its own proof is not part of this item.\par\medskip
\ifdefined\hypertarget\hypertarget{sourcepointer-4}{}\fi
\noindent\textit{Source reference.} Original lemma 0DRX: sources/stacks-chow/tags.txt, source lines 17237--17237. Cross-chapter standard prerequisite cited by the author inside the counted span; it is used as a pointer and its own proof is not part of this item.\par\medskip
\ifdefined\hypertarget\hypertarget{sourcepointer-5}{}\fi
\noindent\textit{Source reference.} Original reference 04HJ: sources/stacks-chow/tags.txt, source lines 5870--5870. Cross-chapter standard prerequisite cited by the author inside the counted span; it is used as a pointer and its own proof is not part of this item.\par\medskip
\ifdefined\hypertarget\hypertarget{sourcepointer-6}{}\fi
\noindent\textit{Source reference.} Original reference 04N4: sources/stacks-chow/tags.txt, source lines 6065--6065. Cross-chapter standard prerequisite cited by the author inside the counted span; it is used as a pointer and its own proof is not part of this item.\par\medskip
\begin{lemma}
\label{lemma-monomorphing-a-component-in-of-the-right-dimension}
Let $\mathcal{T} \hookrightarrow \mathcal{X}$ be a locally
of finite type monomorphism of algebraic stacks,
with $\mathcal{X}$ (and thus also $\mathcal{T}$)
being Jacobson, pseudo-catenary, and locally Noetherian.
Suppose further that $\mathcal{T}$ is irreducible
of some (finite) dimension $d$, and that $\mathcal{X}$ is reduced
and of dimension less
than or equal to $d$.
Then there is a non-empty open substack $\mathcal{V}$ of $\mathcal{T}$ such
that the induced
monomorphism $\mathcal{V} \hookrightarrow \mathcal{X}$ is an open immersion
which identifies
$\mathcal{V}$ with an open subset of an irreducible component of $\mathcal{X}$.
\end{lemma}

\begin{proof}
Choose a smooth surjective morphism $f:U \to \mathcal{X}$ with source a scheme,
necessarily reduced since $\mathcal{X}$ is,
and write $U' = \mathcal{T}\times_{\mathcal{X}} U$.  The base-changed morphism
$U' \to U$ is a monomorphism of algebraic spaces, locally of finite
type, and thus representable
Morphisms of Spaces, Lemma
\ref{spaces-morphisms-lemma-locally-quasi-finite-separated-representable} and
\ref{spaces-morphisms-lemma-monomorphism-loc-finite-type-loc-quasi-finite};
since $U$ is a scheme, so is $U'$.
The projection $f': U' \to \mathcal{T}$ is again a smooth surjection.
Let $u' \in U'$, with image $u \in U$.
Lemma \ref{lemma-base-change-invariance-of-relative-dimension}
shows that $\dim_{u'}(U'_{f(u')}) = \dim_u(U_{f(u)}),$
while $\dim_{f'(u')}(\mathcal{T}) =d
\geq \dim_{f(u)}(\mathcal{X})$ by
Lemma \ref{lemma-irreducible-implies-equidimensional}
and our assumptions on $\mathcal{T}$ and $\mathcal{X}$.
Thus we see that
\begin{equation}
\label{equation-dim-inequality}
\dim_{u'} (U') = \dim_{u'} (U'_{f(u')}) + \dim_{f'(u')}(\mathcal{T})
\\
\geq \dim_u (U_{f(u)}) + \dim_{f(u)}(\mathcal{X}) = \dim_u (U).
\end{equation}
Since $U' \to U$ is a monomorphism, locally of finite type,
it is in particular unramified,
and so by the \'etale local structure of unramified morphisms
\'Etale Morphisms, Lemma \ref{etale-lemma-finite-unramified-etale-local},
we may find a commutative diagram
$$
\xymatrix{
V' \ar[r]\ar[d] & V \ar[d] \\
U' \ar[r] & U
}
$$
in which the scheme $V'$ is non-empty,
the vertical arrows are \'etale,
and the upper horizontal arrow is a closed immersion.
Replacing $V$ by a quasi-compact open subset
whose image has non-empty intersection with the image of $U'$,
and replacing $V'$ by the preimage of $V$, we may further
assume that $V$ (and thus $V'$) is quasi-compact.
Since $V$ is also locally Noetherian,
it is thus Noetherian, and so is the union of finitely many irreducible
components.

\medskip\noindent
Since \'etale morphisms preserve pointwise dimension
Descent, Lemma \ref{descent-lemma-dimension-at-point-local}
we deduce from (\ref{equation-dim-inequality})
that for any point $v' \in V'$,
with image $v \in V$, we have
$\dim_{v'}( V') \geq \dim_v(V)$.
In particular, the image of $V'$ can't be contained in the intersection
of two distinct irreducible components of $V$, and so we may find
at least one irreducible open subset of $V$ which has non-empty intersection
with $V'$; replacing $V$ by this subset, we may assume that $V$ is integral
(being both reduced and irreducible).  From the preceding inequality
on dimensions, we conclude that the closed immersion $V' \hookrightarrow V$
is in fact an isomorphism.
If we let $W$ denote the image of $V'$
in $U'$, then $W$ is a non-empty
open subset of $U'$ (as \'etale morphisms are open),
and the induced monomorphism $W \to U$ is \'etale
(since it is so \'etale locally on the source, i.e.\ after pulling back
to $V'$), and hence is an open immersion (being an \'etale monomorphism).
Thus, if we let $\mathcal{V}$ denote the image of $W$ in $\mathcal{T}$,
then $\mathcal{V}$ is a dense (equivalently, non-empty) open substack of
$\mathcal{T}$,
whose image is dense in an irreducible component of $\mathcal{X}$.
Finally,
we note that the morphism is $\mathcal{V} \to \mathcal{X}$ is smooth
(since its composite
with the smooth morphism $W\to \mathcal{V}$ is smooth),
and also a monomorphism, and thus is an open immersion.
\end{proof}

\end{document}
